\chapter{Platone di Atene}
\label{cap:platone:profilo}

\section{Le fonti}

\subsection{Il corpus platonico e la sua trasmissione}

Di Platone possediamo un corpus di scritti eccezionalmente ampio rispetto a quello conservato dei filosofi greci dei secoli precedenti. La tradizione ha tramandato sotto il suo nome trentaquattro dialoghi, l'\emph{Apologia di Socrate} e tredici lettere. Non tutte le opere attribuitegli sono considerate autentiche, ma l'attribuzione dei dialoghi maggiori non presenta generalmente difficoltà paragonabili a quelle degli scritti minori e del corpus epistolare.\footnote{Per il quadro del corpus e i problemi di attribuzione, cfr.\ \cite[95--96]{palumbo2015}; Diogene Laerzio, \emph{Vite dei filosofi}, III, 50--62.}

La tradizione antica organizzò gli scritti secondo diversi ordinamenti. Quello in nove tetralogie, attribuito al grammatico Trasillo e documentato da Diogene Laerzio (\emph{Vite dei filosofi}, III, 56--62), ebbe particolare fortuna. Tale ordinamento non corrisponde necessariamente alla successione cronologica delle opere.

I dialoghi costituiscono la fonte fondamentale per lo studio della filosofia platonica. La loro forma letteraria richiede tuttavia un'attenzione particolare: Platone rappresenta discussioni nelle quali sono i personaggi, non l'autore, a formulare le tesi e a sviluppare le argomentazioni.\footnote{Cfr.\ \cite[95--98]{palumbo2015}.}

\subsection{Testimonianze antiche e problemi di attendibilità}

Le principali testimonianze esterne provengono da Aristotele, che frequentò l'Accademia, e dalla tradizione biografica e dossografica raccolta da Diogene Laerzio.

Aristotele costituisce una fonte privilegiata per le dottrine attribuite a Platone e per l'ambiente accademico: si vedano, per esempio, \emph{Metafisica}, I, 6, 987a29--987b14. Le sue ricostruzioni, tuttavia, appartengono a un'indagine filosofica autonoma e richiedono di distinguere la posizione riferita dalla critica che ne viene proposta.

Le notizie biografiche presentano difficoltà differenti. Diogene Laerzio, vissuto molti secoli dopo Platone, conserva informazioni tratte da autori precedenti, ma anche genealogie e aneddoti. Una fonte particolare è la \emph{Lettera VII}, che contiene un racconto in prima persona delle vicende politiche siciliane, ma la cui autenticità rimane controversa.\footnote{Per la discussione moderna dell'autenticità, cfr.\ \cite{burnyeatfrede2015}; per l'utilizzo biografico della lettera, cfr.\ \cite[94--95]{palumbo2015}.}

La documentazione consente dunque lo studio diretto dei problemi filosofici rappresentati nei dialoghi principali; molti particolari della vita dell'autore richiedono invece una valutazione critica delle testimonianze.

\section{La vita}

\subsection{Atene, Socrate e il contesto politico}

\subsubsection{Origine familiare e cronologia}

Platone apparteneva a una famiglia ateniese di elevata condizione sociale. La tradizione ne indica come genitori Aristone e Perictione e collega la famiglia materna a Solone e quella paterna all'antica dinastia regale ateniese. Queste genealogie, raccolte da Diogene Laerzio (\emph{Vite dei filosofi}, III, 1), non sono verificabili in ogni particolare e comprendono elementi celebrativi.\footnote{Cfr.\ \cite[94]{palumbo2015}.}

Nel corpus platonico compare una notizia relativa alla famiglia.

\begin{testimonianza}{Platone, \emph{Apologia di Socrate}, 34a}{test:platone:adimanto}
\noindent{\upshape\small\raggedright\textbf{Originale greco.} Ἀδείμαντος, ὁ Ἀρίστωνος, οὗ ἀδελφὸς οὑτοσὶ Πλάτων.\par}
\smallskip
\noindent{\upshape\small\textbf{Resa italiana.}}\par
Adimanto, figlio di Aristone, del quale questo Platone è fratello.
\end{testimonianza}

\noindent{\footnotesize\emph{Traduzione italiana redazionale; testo di controllo: testo greco dell'\emph{Apologia}, 34a, riproduzione digitale della Biblioteca della lingua greca (\texttt{greek-language.gr}).}}

\noindent\textbf{Analisi critica.}
Nel discorso attribuito a Socrate, Adimanto è indicato come figlio di Aristone e fratello di Platone. La notizia concorda con la genealogia riferita da Diogene Laerzio, ma appartiene a una rappresentazione letteraria del processo; non documenta indipendentemente la presenza di Platone in tribunale.

La cronologia tradizionale colloca la nascita di Platone nel 428/427 a.C. sulla base delle testimonianze antiche. Una diversa ricostruzione proposta da Robin Waterfield indica il 424/423 a.C.\footnote{Diogene Laerzio, \emph{Vite dei filosofi}, III, 2--3; cfr.\ \cite[cap.~1]{waterfield2023}.} Anche il luogo di nascita presenta qualche incertezza: Diogene Laerzio (III, 3) conserva una tradizione che lo colloca a Egina, senza mettere in discussione l'appartenenza di Platone alla cittadinanza ateniese.

Secondo un'altra notizia tramandata da Diogene Laerzio (III, 4), il suo nome originario sarebbe stato Aristocle. La tradizione è discussa: Waterfield ritiene che si sia sempre chiamato Platone.\footnote{Cfr.\ \cite[cap.~1]{waterfield2023}; per la lettura del nome in Palumbo e il rinvio a D.\ Sedley, cfr.\ \cite[94]{palumbo2015}.} Non è dunque possibile presentare \emph{Aristocle} come nome di nascita accertato.

\subsubsection{La formazione e Socrate}

Platone ricevette probabilmente l'educazione propria della sua condizione sociale, comprendente la formazione letteraria, la musica e la ginnastica. Le testimonianze sulla giovinezza sono prevalentemente posteriori e non consentono di ricostruire con precisione il percorso formativo.\footnote{Cfr.\ \cite[94]{palumbo2015}; Diogene Laerzio, \emph{Vite dei filosofi}, III, 4--5.}

Il rapporto con Socrate costituisce invece un dato fondamentale, documentato dalla produzione dialogica e dalla tradizione antica. La data dell'inizio della loro frequentazione, tradizionalmente collocata intorno al 408 a.C., non è accertabile con precisione (Diogene Laerzio, III, 6).

La giovinezza di Platone coincise con una fase travagliata della storia ateniese. La sconfitta nella guerra del Peloponneso, nel 404 a.C., fu seguita dal regime oligarchico dei Trenta e dal successivo ripristino della democrazia. Nel 399 a.C. Socrate fu processato e condannato a morte. L'avvenimento occupa una posizione centrale nell'\emph{Apologia}, nel \emph{Critone} e nel \emph{Fedone}.

Il significato politico attribuito a queste esperienze è ricostruito soprattutto attraverso la \emph{Lettera VII} (324b--326b): l'autore descrive la propria delusione nei confronti delle istituzioni ateniesi, oligarchiche e democratiche, e collega la possibilità di una politica giusta alla conoscenza filosofica. Il valore autobiografico del racconto dipende però dall'autenticità della lettera.\footnote{Cfr.\ \cite[94--95]{palumbo2015}; sulla controversia di autenticità, \cite{burnyeatfrede2015}.}

\subsection{I viaggi in Sicilia e la fondazione dell'Accademia}

\subsubsection{I soggiorni a Siracusa}

La tradizione colloca tre soggiorni di Platone a Siracusa, generalmente datati intorno al 388/387, al 367 e al 361 a.C. I rapporti con la corte siracusana e con Dione sono riferiti dalla \emph{Lettera VII}, da Diogene Laerzio e da Plutarco (\emph{Vita di Dione}); le circostanze e le motivazioni dei singoli viaggi non sono ricostruibili con eguale certezza.\footnote{Cfr.\ \cite[94--95]{palumbo2015}; Diogene Laerzio, \emph{Vite dei filosofi}, III, 18--23; Plutarco, \emph{Vita di Dione}, 4--5, 11--18.}

Nel primo soggiorno Platone avrebbe conosciuto Dione, legato alla dinastia di Dionisio il Vecchio. Dopo la morte del tiranno, Dione cercò di coinvolgerlo nell'educazione politica del successore, Dionisio il Giovane.

La \emph{Lettera VII} presenta il ritorno a Siracusa come un tentativo di tradurre in azione una concezione filosofica del governo.

\begin{testimonianza}{\emph{Lettera VII}, 328b--c (attribuita a Platone)}{test:platone:lettera-vii}
\noindent{\upshape\small\raggedright\textbf{Originale greco.} ὅθεν μοι σκοπουμένῳ καὶ διστάζοντι πότερον εἴη πορευτέον καὶ ὑπακουστέον ἢ πῶς, ὅμως ἔρρεψε δεῖν, εἴ ποτέ τις τὰ διανοηθέντα περὶ νόμων τε καὶ πολιτείας ἀποτελεῖν ἐγχειρήσοι, καὶ νῦν πειρατέον εἶναι· πείσας γὰρ ἕνα μόνον ἱκανῶς πάντα ἐξειργασμένος ἐσοίμην ἀγαθά.\par}
\smallskip
\noindent{\upshape\small\textbf{Resa italiana.}}\par
Mentre riflettevo, incerto se dovessi partire e accogliere l'invito, prevalse tuttavia la convinzione che, se mai avessi tentato di realizzare ciò che avevo concepito riguardo alle leggi e alla costituzione, fosse quello il momento di provarci: persuadendo adeguatamente un solo uomo, avrei potuto realizzare ogni bene.
\end{testimonianza}

\noindent{\footnotesize\emph{Traduzione italiana redazionale; testo di controllo: testo greco della \emph{Lettera VII}, 328b--c, riproduzione digitale ScholaThomistica.}}

\noindent\textbf{Analisi critica.}
Il passo stabilisce un rapporto esplicito fra elaborazione filosofica ed esercizio del governo. L'autore presenta l'intervento politico come un tentativo di realizzare concezioni precedentemente elaborate sulle leggi e sulla costituzione.

La testimonianza non è automaticamente autobiografica. Mentre alcuni studiosi accettano l'autenticità della lettera, Myles Burnyeat e Michael Frede sostengono che sia un'opera pseudoepigrafa.\footnote{Cfr.\ \cite{burnyeatfrede2015}; per il quadro biografico, \cite{waterfield2023}.} La controversia riguarda quindi anche il valore storico delle motivazioni attribuite a Platone.

I tentativi siracusani non realizzarono la trasformazione del governo auspicata nella tradizione platonica. Le tensioni fra Dione e Dionisio il Giovane condussero a conflitti culminati nella presa del potere da parte di Dione nel 357 a.C. e nel suo assassinio nel 354 a.C.\footnote{Cfr.\ Plutarco, \emph{Vita di Dione}, 47--57; \cite[95]{palumbo2015}.}

\subsubsection{L'Accademia}

La fondazione dell'Accademia è tradizionalmente collocata intorno al 387 a.C., dopo il primo soggiorno siciliano. La data è approssimativa: non conosciamo un atto formale di fondazione e la comunità filosofica potrebbe essersi costituita gradualmente.\footnote{Cfr.\ \cite[95]{palumbo2015}; \cite[cap.~5]{waterfield2023}.}

Il nome deriva dall'area dell'Accademia, fuori dalle mura di Atene, presso un ginnasio. Vi operarono filosofi e studiosi interessati anche alla matematica, all'astronomia e alla politica. Non era un'università nel senso moderno, e la sua organizzazione concreta rimane in parte incerta.

Fra coloro che la frequentarono figurano Speusippo, Senocrate, Eudosso e Aristotele. Quest'ultimo vi trascorse circa vent'anni. Platone morì nel 348/347 a.C., secondo la cronologia generalmente accolta; alla direzione della scuola gli succedette Speusippo.\footnote{Per la successione accademica, Diogene Laerzio, \emph{Vite dei filosofi}, IV, 1--3; per la cronologia biografica, III, 2--3.}

\section{Le opere}

\subsection{Dialoghi e lettere: attribuzioni e autenticità}

Quasi tutta la produzione filosofica conservata di Platone è costituita da dialoghi. La loro peculiarità non consiste soltanto nell'alternanza di domande e risposte, ma nella rappresentazione di un'indagine nella quale gli interlocutori formulano ipotesi, argomentano e sottopongono a esame le rispettive posizioni.\footnote{Cfr.\ \cite[95--98]{palumbo2015}.}

I personaggi sono spesso figure storiche. Socrate occupa una posizione dominante in numerose opere, ma non deve essere identificato semplicemente con il Socrate storico: la distinzione fra il maestro e l'elaborazione dell'allievo è un problema interpretativo fondamentale. Il corpus comprende inoltre lettere tramandate sotto il nome di Platone, il cui valore documentario dipende dalle questioni di autenticità già ricordate.

\subsection{La cronologia controversa delle opere}

La successione cronologica degli scritti non è conosciuta con certezza. La ricerca moderna ha utilizzato i riferimenti interni e l'analisi stilometrica, cioè lo studio delle caratteristiche ricorrenti della lingua e dello stile.\footnote{Cfr.\ \cite[106--108]{palumbo2015}.}

Una classificazione largamente impiegata distingue dialoghi giovanili, della maturità e tardi. Ai primi vengono comunemente ricondotti \emph{Eutifrone}, \emph{Lachete} e \emph{Carmide}; fra quelli della maturità si collocano abitualmente \emph{Fedone}, \emph{Simposio}, \emph{Fedro} e \emph{Repubblica}. Un gruppo stilisticamente più omogeneo comprende opere considerate tarde, fra le quali \emph{Sofista}, \emph{Politico}, \emph{Filebo}, \emph{Timeo} e \emph{Leggi}.

Questa classificazione non equivale a una cronologia assoluta: la collocazione di alcuni dialoghi rimane discussa e la somiglianza stilistica non basta a datarli con precisione. Non deve inoltre essere confusa con una ricostruzione necessariamente evolutiva del pensiero: differenze dottrinali possono dipendere anche dai problemi e dalle situazioni argomentative.\footnote{Cfr.\ \cite[106--110]{palumbo2015}.}

\subsection{La forma dialogica e la voce dell'autore}

La scrittura platonica comprende discussioni brevi, lunghi discorsi, racconti mitici e indagini che non raggiungono una conclusione definitiva. La forma dialogica non autorizza a negare che l'autore abbia convinzioni filosofiche, ma impone di ricostruirle attraverso l'argomentazione complessiva, non isolando affermazioni pronunciate dai personaggi.

Il rapporto fra dialoghi scritti e insegnamento orale è particolarmente discusso. Nel \emph{Fedro} (274b--278e) Socrate esamina i limiti della scrittura; la \emph{Lettera VII} (341b--342a), se autentica, offre un'ulteriore testimonianza sulle difficoltà di comunicare il sapere filosofico mediante uno scritto. Il tema sarà approfondito nel capitolo~\ref{cap:platone:dialogo}.\footnote{Cfr.\ \cite[108--110]{palumbo2015}.}

\section{Il pensiero}

\subsection{La ricerca filosofica e il problema della definizione}

L'eredità socratica costituisce uno dei punti di partenza della riflessione platonica. Nei dialoghi la filosofia si presenta anzitutto come ricerca che sottopone a esame le opinioni e richiede agli interlocutori di giustificare razionalmente le loro affermazioni. Nel \emph{Gorgia} tale concezione emerge dal confronto con la retorica.

\begin{testimonianza}{Platone, \emph{Gorgia}, 458a}{test:platone:gorgia-confutazione}
\noindent{\upshape\small\raggedright\textbf{Originale greco.} ἐγὼ δὲ τίνων εἰμί; τῶν ἡδέως μὲν ἂν ἐλεγχθέντων εἴ τι μὴ ἀληθὲς λέγω, ἡδέως δʼ ἂν ἐλεγξάντων εἴ τίς τι μὴ ἀληθὲς λέγοι, οὐκ ἀηδέστερον μεντἂν ἐλεγχθέντων ἢ ἐλεγξάντων· μεῖζον γὰρ αὐτὸ ἀγαθὸν ἡγοῦμαι, ὅσῳπερ μεῖζον ἀγαθόν ἐστιν αὐτὸν ἀπαλλαγῆναι κακοῦ τοῦ μεγίστου ἢ ἄλλον ἀπαλλάξαι.\par}
\smallskip
\noindent{\upshape\small\textbf{Resa italiana.}}\par
E io a quali uomini appartengo? A quelli che accetterebbero volentieri di essere confutati, qualora dicessero qualcosa di falso, e che altrettanto volentieri confuterebbero chi dicesse qualcosa di falso; anzi, a quelli per i quali essere confutati non sarebbe meno piacevole che confutare. Ritengo infatti che sia un bene maggiore liberare se stessi dal male più grande piuttosto che liberarne un altro.
\end{testimonianza}

\noindent{\footnotesize\emph{Traduzione italiana redazionale; testo di controllo: testo greco del \emph{Gorgia}, 458a, riproduzione digitale ScholaThomistica.}}

\noindent\textbf{Analisi critica.}
Socrate attribuisce alla confutazione (\emph{élenchos}, ἔλεγχος) un valore conoscitivo: la discussione deve condurre alla correzione delle convinzioni false, non alla vittoria sull'interlocutore. Il passo presenta una concezione della filosofia come ricerca della verità, distinta dal successo persuasivo nella disputa.\footnote{Cfr.\ \cite[97--98]{palumbo2015}.}

La domanda socratica su che cosa siano la giustizia, il coraggio o la virtù esige una definizione, non una semplice serie di esempi. Questo problema orienta diverse indagini platoniche e introduce la questione delle condizioni stesse del sapere.

\subsection{Conoscenza, idee e dialettica}

La ricerca delle definizioni conduce a interrogarsi sul rapporto fra la molteplicità delle cose e ciò che esse hanno in comune. In alcuni dialoghi emerge la distinzione fra le cose sensibili e le realtà intelligibili, designate con i termini \emph{idéa} (ἰδέα) ed \emph{eîdos} (εἶδος).

\begin{testimonianza}{Platone, \emph{Fedone}, 100c}{test:platone:fedone-bello}
\noindent{\upshape\small\raggedright\textbf{Originale greco.} φαίνεται γάρ μοι, εἴ τί ἐστιν ἄλλο καλὸν πλὴν αὐτὸ τὸ καλόν, οὐδὲ δι᾽ ἓν ἄλλο καλὸν εἶναι ἢ διότι μετέχει ἐκείνου τοῦ καλοῦ· καὶ πάντα δὴ οὕτως λέγω.\par}
\smallskip
\noindent{\upshape\small\textbf{Resa italiana.}}\par
Mi sembra infatti che, se esiste qualcosa di bello oltre al Bello in sé, esso non sia bello per nessun'altra ragione se non perché partecipa di quel Bello; e sostengo lo stesso riguardo a tutte le altre cose.
\end{testimonianza}

\noindent{\footnotesize\emph{Traduzione italiana redazionale; testo di controllo: testo greco del \emph{Fedone}, 100c, riproduzione digitale ScholaThomistica.}}

\noindent\textbf{Analisi critica.}
Socrate distingue le cose particolari dotate di una proprietà dalla realtà intelligibile alla quale esse partecipano. Nel contesto del \emph{Fedone} (100b--d) le forme sono introdotte come ipotesi esplicativa: il passo non è una dimostrazione autonoma della loro esistenza.

La teoria delle idee costituisce un nucleo importante della filosofia dei dialoghi, ma viene sottoposta a difficoltà e obiezioni, soprattutto nel \emph{Parmenide} (130a--135d). Severino interpreta questa dottrina come un approfondimento della ricerca socratica dell'universale e del problema, ereditato dalla tradizione eleatica, del rapporto fra mutamento e immutabilità.\footnote{Cfr.\ \cite[121--126]{severino2010}. Si tratta della ricostruzione di Severino, non di un'affermazione testuale attribuibile indiscriminatamente a Platone.}

Altri dialoghi sviluppano questioni propriamente conoscitive: il \emph{Menone} (80d--86c) esamina la possibilità di apprendere; il \emph{Teeteto} (151d--210d) discute diverse definizioni del sapere; la \emph{Repubblica} (VI, 509d--511e) distingue differenti livelli della conoscenza. Gli argomenti saranno approfonditi nei capitoli dedicati alle idee e alla dialettica.

\subsection{Anima, giustizia e vita politica}

Le questioni conoscitive sono connesse alla riflessione sull'anima (\emph{psychḗ}, ψυχή). Nel \emph{Fedone} l'anima è esaminata in rapporto alla morte e alla conoscenza delle realtà intelligibili; nella \emph{Repubblica} (IV, 436a--441c) la distinzione fra le componenti razionale, irascibile e desiderativa permette di affrontare il problema dei conflitti interiori.

La riflessione sull'anima assume quindi un significato morale: comprendere l'uomo significa interrogarsi sulle condizioni di una vita giusta. Nella \emph{Repubblica} l'indagine sulla giustizia conduce all'esame della città e della formazione dei governanti. Socrate sostiene che i mali delle città non avranno fine finché il potere politico e la filosofia non coincideranno (\emph{Repubblica}, V, 473c--d).

La tesi pone un problema ulteriore, non risolvibile mediante la sola formula del «filosofo-re»: quale sapere rende qualcuno capace di governare giustamente? Il tema sarà sviluppato nei capitoli sulla \emph{Repubblica}.

\subsection{Natura, cosmo e linguaggio}

L'indagine platonica si estende ai problemi della natura, dell'essere e del linguaggio. Nel \emph{Timeo} (27d--29d) la spiegazione dell'ordine cosmico prende avvio dalla distinzione fra ciò che è sempre e ciò che diviene, con il problema dello statuto conoscitivo del discorso sulla natura.

Nel \emph{Sofista} (237a--264b) la ricerca di una definizione del sofista conduce a un esame dell'essere, del non essere e della possibilità del discorso falso. La varietà delle indagini non autorizza, da sola, a ricostruire un sistema unitario: i loro rapporti devono emergere dallo studio delle argomentazioni.

\section{L'eredità}

\subsection{L'Accademia e la tradizione successiva}

L'influenza di Platone si esercitò inizialmente attraverso l'Accademia. Speusippo e Senocrate svilupparono orientamenti differenti, mostrando che la continuità della scuola non comportava l'adesione immutabile alle medesime dottrine. Aristotele riprese numerosi problemi platonici, criticandone alcune soluzioni.\footnote{Per la successione della scuola, Diogene Laerzio, \emph{Vite dei filosofi}, IV, 1--15; per il rapporto filosofico con Platone, Aristotele, \emph{Metafisica}, I, 6 e XIII, 4.}

La tradizione accademica e platonica attraversò orientamenti diversi, dall'Accademia ellenistica alle forme di platonismo sviluppatesi in età imperiale e tardoantica. La storia di tali interpretazioni non può essere identificata semplicemente con il pensiero del fondatore.

\subsection{La pluralità delle interpretazioni di Platone}

La ricostruzione della filosofia platonica rimane oggetto di un dibattito storiografico: una linea interpretativa attribuisce particolare valore all'insegnamento orale dell'Accademia e cerca di riconoscere i principi di un sistema unitario; un'altra insiste sul carattere aperto e problematico della ricerca rappresentata nei dialoghi.

Palumbo presenta queste prospettive e richiama, attraverso Mario Vegetti, una posizione intermedia: individuare alcuni nuclei teorici relativamente costanti senza trasformare necessariamente la filosofia platonica in un sistema definitivamente concluso.\footnote{Cfr.\ \cite[109--110]{palumbo2015}.}

Lo studio di Platone richiede pertanto di distinguere il testo dei dialoghi, le testimonianze antiche e le ricostruzioni degli interpreti. L'unità del suo pensiero non deve essere assunta come premessa indiscutibile: costituisce essa stessa un problema della ricerca storico-filosofica.

% Il repertorio comprende soltanto studi moderni effettivamente citati.
% I rinvii ai passi antichi rimangono nel testo nelle forme canoniche.
\addsec{Bibliografia}
\printbibliography[heading=none]
